\documentclass[11pt]{article}
\usepackage{raeez-math-template}
\numberwithin{equation}{section}
\title{Meridian Tori in a Two-Filling Attachment}
\author{Raeez Lorgat}
\date{}
\hypersetup{pdfsubject={Integral second homology and marked meridian tori}}
\begin{document}
\maketitle
\begin{abstract}
Two finite torus fillings determine a compact real six-manifold with a
torus-bundle boundary. We compute its integral second homology and the
images of two marked meridian tori. A double mapping cylinder retains the
attaching paths and reduces the calculation to explicit core surfaces.
The two boundary classes form an integral basis. Attaching a cusp whose
specialization has these two classes as its kernel therefore kills the
second homology.
\end{abstract}
\input{meridian-tori}
\end{document}
